% VestigeKV defense document: anticipated reviewer questions with the full
% detail that the whitepaper compresses out. Companion to
% docs/whitepaper/whitepaper.tex; every answer cites the whitepaper section
% it expands.
\documentclass[11pt]{article}
\usepackage[margin=1in]{geometry}
\usepackage{enumitem}
\usepackage{hyperref}

\title{VestigeKV --- Defense Document\\{\large Anticipated questions, with the
detail the whitepaper compresses out}}
\date{}

\begin{document}
\maketitle

\noindent Each section below is one anticipated question. The whitepaper
carries the compressed statement; this document carries the full argument
and the measured numbers behind it.

%======================================================================
\section{Q1: At a 99.8\% trigger rate, is tier 1 still necessary at long
context? Does recall depend on it?}

Yes, and recall is structurally defined in terms of it.

\textbf{Trigger frequency is not workload share.} At $S{=}512$k a step
attends $\approx\!16{,}384$ tier-1 rows and fetches $\approx\!120$ recalled
ones ($0.024\%$ of the archive): tier 1 supplies over $99\%$ of the attended
rows and essentially all of the attention mass; recall is the patch, not the
payload. The correct mental model is a cache hierarchy --- tier 1 is the
\emph{hit path} (query-independent, prefetched at prefill, riding the fixed
in-graph channel at zero per-step cost), recall is the \emph{miss path}
(query-dependent, riding the variable channel). ``At least one miss on
$99.8\%$ of steps, at $0.024\%$ miss volume'' describes a healthy cache, not
a redundant one. Removing tier 1 forces every relevant row through the
variable scan-and-fetch path every step --- the fired set balloons from
${\sim}120$ rows to the entire important set, converting cheap fixed reads
into scanned fetches in every regime.

\textbf{Recall depends on tier 1 three times over.} (a)~The fire criterion
is $\mathrm{score}_a > \mathrm{max}_1$, and $\mathrm{max}_1$ \emph{is} tier
1's content --- without the kept set the competition standard does not
exist. (b)~The entropy gate is computed over the kept-tier softmax.
(c)~Calibration defines a hard query as one whose target is \emph{archived},
i.e.\ relative to the tier-1 partition. The counterfactual was run and
rejected in the pre-registered record: trigger designs without the kept-max
competition either recalled poorly (recall@4 $=0.489$) or fired on $99.5\%$
of rows --- it is precisely the competition against tier 1's optimum that
tames firing from $99.5\%$ to $0.024\%$. Remove tier 1 and the mechanism
collapses back to the similarity top-$k$ shape this design supersedes.

%======================================================================
\section{Q2: How does this differ from ArkVale?}

ArkVale (Chen et al., 2024) is the closest prior design --- evict, summarize,
recall on demand --- and a port team should know exactly where this tier
departs from it:
\begin{enumerate}[itemsep=1pt]
\item \textbf{Unit and summary.} ArkVale summarizes evicted \emph{pages}
with lossy bounding digests and recalls whole pages. This tier indexes
individual \emph{rows} with an exact 64-dim branch summand, a rank-$r$
content sketch, and a scalar residual certificate; recalled rows are
bit-exact (no lossy digest exists anywhere, per the no-quantization rule).
\item \textbf{Why digests cannot be tight under RoPE.} A page digest must
cover its keys; rotation spreads a page's keys across phases, so the cover
is structurally loose --- measured digest layer-uniformity $0.672$ (NoPE)
vs.\ $0.024$ (RoPE) in the reference record. ArkVale operates in the loose
regime; this tier is built for the NoPE regime where per-row summaries are
meaningful at all.
\item \textbf{Certified competition vs.\ similarity scoring.} ArkVale admits
pages by digest-similarity ranking. Here a row fires only when its certified
score beats the query's own kept-tier maximum, with $z$ self-calibrated to a
90\% recall target from the prefix --- an admission \emph{criterion with a
guarantee}, not a similarity heuristic.
\item \textbf{Eviction signal.} ArkVale's eviction relies on query-history
importance; tier-1 eviction here is query-independent ($\sigma$ over the
branch), so the whole partition exists before any query does --- the setting
ArkVale cannot serve.
\item \textbf{Scan economy (design intent, with a measured caveat).} The
entropy gate is meant to skip whole scans when the kept tier is already
confident, where ArkVale scans per step. On the \texttt{sglang} port the
gate is \emph{off} even after calibration is repaired --- with $3$--$63$
hard queries per layer the fitted threshold would pass more than $60\%$ of
queries, so the self-disable rule correctly declines to install a filter
that filters nothing. The conformal $z$, which shares the same hard set,
does calibrate. Budget a full scan per head per step; the gate is an
optimization the port does not currently collect.
\end{enumerate}

\paragraph{NoPE-specific optimizations in the recall path.} Each departure
is enabled by a specific NoPE property, with a RoPE counterfactual
explaining why ArkVale cannot adopt it:
(i)~\emph{build-once, never-stale row index} --- scores are fixed bilinear
forms of row content, so per-row entries computed at the compression event
stay valid for the whole decode; under RoPE the effective key drifts with
query--key distance, which is exactly why ArkVale must fall back to loose
per-page bounding digests.
(ii)~\emph{a free exact summand} --- the never-rotated 64-dim branch makes
$q_{\mathrm{rope}}\cdot\mathrm{side}_a$ an exact component of the true score
at 11\% of the row read; rotation leaves no exact sub-component to
precompute.
(iii)~\emph{a static certificate} --- the sketch-truncation error is a
per-row constant $\varrho_a$ only because the content form is fixed; a
rotating subspace admits no static bound, hence similarity heuristics
instead of certified admission.
(iv)~\emph{one-shot calibration} --- $z$ and the gate threshold are
calibrated once from the prefix and remain valid because score distributions
are position-stationary; under RoPE prefix calibration drifts as the rotary
bands sweep.
(v)~\emph{zero positional fix-up on recall} --- fetched rows splice into the
same softmax with no re-rotation, which is what lets the stock kernel serve
the merged partition unmodified.

\paragraph{Memory placement: no paging.} Eviction here is a partition of
\emph{indices}, never a movement of data: in the standard (GPU-resident)
configuration every row --- kept and archived alike --- stays in place in
the one GPU pool, and a recall adds \emph{indices} to the step's attended
set; the kernel gathers the rows where they lie. There is no page, no
resident set, no replacement, and therefore no thrashing, by construction
(the compression benefit in this configuration is bandwidth and compute, not
VRAM). The host-offload variant does move archived rows to host memory once,
at the compression event, off the critical path; a fire then costs a
one-way, row-granular fetch (rows are immutable: no dirty state, no
write-back, no displacement of anything resident), bounded under the cap by
$\mathrm{topj}\times H$ rows per step per layer. Neither configuration
performs swap-in/swap-out in ArkVale's sense.

\paragraph{The degeneration duality.} The two designs fail in opposite
directions, and the contrast is the sharpest way to state the difference. A
fixed page budget never inflates per-step attention --- ArkVale does
\emph{not} degenerate to naive MLA on recall --- but under pressure it
degrades \emph{quality}: pages beyond the budget are simply absent (plus
transfer thrashing when the relevant set churns), and the only cure is
inflating the budget toward the full cache, which un-does the compression.
The certified competition degrades in the dual direction: uncapped, a
pathological step can attend everything (naive MLA plus scan) while quality
never degrades --- the fetched rows are exact. The explicit cap is what
turns this into an operator choice: quota designs hard-wire the degradation
onto the quality side silently; here the side and the bound are stated.

%======================================================================
\section{Q3: What is the validation protocol?}

\begin{enumerate}[itemsep=2pt]
\item \textbf{Kill-switch parity}: with the wrapper disabled, outputs must
be call-for-call identical to stock. This isolates plumbing bugs from
algorithm effects.
\item \textbf{Full-arm parity}: with selection forced to keep \emph{all}
rows, per-step outputs must match stock bit-for-bit. (The reference A/B
harness switches arms on a live server this way; both arms then share one
graph and code path, which is also the honest benchmark protocol.)
\item \textbf{Selector equivalence}: $\sigma$ and the keep set must match
the naive reference implementation --- bit-exactly for the fallback path,
decision-exactly (identical keep masks) for the fused kernels
(\texttt{engine/test/manual/test\_vestigekv\_equiv.py}, self-contained).
\item \textbf{Quality gates} (uncapped, $\rho{=}1/32$): needle retrieval
1.00 at $8\times$--$32\times$ / 8k--65k; teacher-forced $\Delta$bpb vs the
uncompressed engine $\le 0.15\%$; generation under the cap bit-identical.
\item \textbf{Recall-tier gates}: with the index built from the prefix (no
training), recovery at $1/128$ compression must reach $1.00$ in the serving
loop at 8k--32k, and the long-context necessity signature must reproduce: a
high per-step trigger rate at $S{\sim}512$k with quality passing (the
reference measures $99.8\%$).
\item \textbf{Scan-precision gates} (the 16-bit index): fire-set stability
--- the fired row set under the 16-bit operands must match the fp32 arm
(Jaccard $\ge 0.99$, top-8 coverage unchanged), each arm calibrated
independently, since the conformal threshold is refit after rounding ---
plus a serving-level greedy + needle A/B.
\item \textbf{Fusion parity}: every fused kernel keeps its pre-fusion
implementation as a bit-parity oracle (fire/fetch sets and CSR contents
identical on shared inputs); end-to-end, the fused pipeline's 2048-token
greedy output must be token-identical to the unfused one.
\item \textbf{Regression tests worth porting} (each caught a real bug):
padded-batch refresh must not overrun the unpadded
$\mathrm{out\_cache\_loc}$; a re-used slot must not inherit the previous
request's kept state; the guard of each test must \emph{fail} on the broken
form (a check must be demonstrated to fail before it is trusted).
\end{enumerate}

%======================================================================
\section{Q4: What defect classes did the port observe, and how were they
mitigated?}

\begin{enumerate}[itemsep=2pt]
\item \textbf{Stale per-slot state on slot reuse.} Runtime request slots are
recycled. Any decode-side cache keyed by slot (the kept tables, tier-2
index) must be invalidated when a new request prefills the slot ---
otherwise the new request attends the \emph{previous} request's rows. This
bug produced ``compression enabled but attention time unchanged'' and was
only caught by a stage-separated kernel profile.
\item \textbf{Selection work inside the decode window.} Building indices
lazily at the first decode step put $B \times L_{\mathrm{MLA}}$ rFFTs into
one step and erased the measured gain at $B{=}16$. All $\sigma$ work belongs
at prefill / block close (the frozen-selection corollary guarantees this is
exact).
\item \textbf{Per-request host loops on the step path.} A Python-level loop
of tiny copies ($\sim$100 ops/step at $B{=}16$) costs milliseconds; the
vectorized pack replaces it with $\sim$12 ops.
\item \textbf{Hidden device sync inside graph capture.} See the
capture-safety rule in the whitepaper's decode-path section.
\item \textbf{Padded-lane overruns.} See the padding trap in the same
section; also give every padded lane a pad row, never an empty range.
\item \textbf{Arch-specific JIT of unrelated ops.} On new hardware,
auxiliary ops (e.g.\ sampling kernels) may fail platform checks and kill
the server at first generate --- pin arch-agnostic fallbacks for everything
outside the attention path before debugging the attention path.
\item \textbf{Graph-replay reality: a backend's \texttt{forward\_*} python
does not run on the replay path.} Only the metadata hooks execute; anything
a per-step algorithm needs must live in the fixed-address buffers those
hooks can touch. Placing a per-step tensor copy inside the forward call
either invalidates the graph (silent fall back to eager) or replays a stale
capture-time value. The reference port measured the former: $138 \to 94$
tok/s from one such copy, with the profiler showing GPU busy $6.4$\,ms
against a $279$\,ms step. Design the per-step data flow around this
constraint before writing the recall path, not after.
\item \textbf{Isolate with a configuration matrix before attributing
anything.} A performance investigation on this port spent ten successive
hypotheses --- each reasoned from the code, each microbenchmarked, each
falsified --- because the comparison itself was never validated first. The
baseline had been measured on one upstream revision and the regression on
another; the two were not comparable at all. Two cheap moves prevent this
class of waste: reproduce the baseline on the \emph{current} tree before
calling anything a regression, and fill a small matrix (stock backend vs.\
wrapper-with-feature-off vs.\ wrapper-with-feature-on) $\times$ (old
revision vs.\ new) so the failing cell names the variable. Record the
upstream commit next to every serving number; never compare across a rebase.
\item \textbf{Host-side cost is invisible to isolated microbenchmarks.}
Five hypotheses (scan GEMV, index pack, bookkeeping, extra attended rows,
device-to-host syncs) were each measured in isolation, each came back cheap,
and each was wrong about the whole: a single-call microbenchmark has no
pipeline to stall. Only a stage-separated trace of the live server
(busy-vs-gap) located the cost. Use microbenchmarks to price kernels, never
to exonerate a host-side path.
\item \textbf{Benchmark hygiene.} Warm each graph bucket before measuring
(first-use capture cost is large); interleave A/B arms on one live server;
never compare across server restarts; measure with the serving framework's
own benchmark tool and state the one exception if any. Two corollaries the
port learned the hard way: (a) the client's inter-token timestamps degrade
under stream batching above ${\sim}240$k context, so the server-side decode
log is the per-token authority there; (b) launch-to-launch network variance
($\pm 15\%$ on the fixed terms) can invert a small cross-launch ratio ---
collect each arm's curve within a single launch and read cross-arm claims as
$\Delta$-vs-4k, and never switch the baseline itself when reducing noise ---
a same-launch full-arm A/B removes launch variance but compares against a
\emph{different} (faster) baseline than the stock backend, which silently
inflates the ratio.
\item \textbf{Second-graph launch tax (fixed, batch-independent).} An
earlier build replayed the recall scan as a \emph{second} CUDA graph per
step; its host-side \texttt{cudaGraphLaunch} waited a full GIL slice and
cost ${\sim}0.9$\,ms/step that did not scale with batch --- enough to sink
bs=1 below $1.0\times$ on its own. The fix is structural, not a knob: bake
the whole recall step into the decode model graph, one launch per step. A
knob-level attempt (shortening the GIL switch interval) helped the target
microcost but regressed the process globally and was reverted; the lesson is
that a per-step fixed cost is worth a structural fix, not a tuning band-aid.
\item \textbf{Capacity-grid launch floor.} The in-graph kernels launch a
grid sized for the worst-case context and batch, so at small live batch the
placeholder pairs cost real launch time even though they compute nothing.
Two mitigations, both content-independent: give each program a whole
1024-row compact bucket (a $16\times$ smaller grid than one-block programs),
and clip the grid per batch-size class to the live pair count so
placeholders are not launched at all. What remains is the platform's own
per-batch-class graph overhead, which the dense arm pays too.
\end{enumerate}

\end{document}
